\chapter{Ziffern und Stellenwertsysteme}
\label{b10guide:ch:numerals}

Dieses Kapitel behandelt natürliche Zahlen und ihre endlichen Zahlzeichen.
Die Indexschreibweise für Ziffern führt keine allgemeine Folgentheorie ein.
Die Begriffe knüpfen an das letzte Kapitel des Hauptbands an;
Existenz, Eindeutigkeit und Ziffernauswertung werden hier als
zentrale Folgerungen ergänzt. Zum Abgleich dient \cite{natdigits}.

\section{Basis, Auswertung und Normalform}

\begin{definition}[Basis, Ziffern und Zahlzeichen]
\label{b10guide:def:digits}
Eine \emph{Basis} ist eine natürliche Zahl $b\ge2$.
Die zugehörigen Ziffern sind die Elemente von
$D_b:=\{d\in\mathbb N\mid d<b\}$.
Für eine Länge $k\ge1$ sind Zifferndaten eine Funktion
\[
 d\colon\{i\in\mathbb N\mid i<k\}\longrightarrow D_b,
 \qquad d_i:=d(i).
\]
Das Zahlzeichen $(d_{k-1}\cdots d_1d_0)_b$ hat den Wert
\[
 V_b(d):=\sum_{i<k}d_i b^i.
\]
Die rechte Ziffer gehört also zur Potenz $b^0=1$.
Für $k=1$ wird nur $(d_0)_b$ geschrieben.
\end{definition}

\begin{definition}[Normalisierte Darstellung]
Eine Darstellung der Länge $k\ge1$ heißt \emph{normalisiert},
wenn $k=1$ oder $d_{k-1}\ne0$ gilt.
Die Darstellung der Null ist das einzelne Zahlzeichen $(0)_b$.
\end{definition}

\begin{theorem}[Existenz und Eindeutigkeit der Stellenwertdarstellung]
\label{b10guide:thm:digits}
Für jede Basis $b\ge2$ und jede Zahl $n\in\mathbb N$ gibt es
genau eine Länge $k\ge1$ und genau eine normalisierte Ziffernbelegung
$d$ dieser Länge mit
\[
 n=\sum_{i<k}d_i b^i.
\]
Für $n>0$ ist die höchste Ziffer $d_{k-1}$ positiv.
Ohne die Normalisierungsbedingung sind zusätzliche führende Nullen möglich.
\end{theorem}

\begin{theorem}[Stellenzahl und Potenzschranken]
Für jede Ziffernbelegung $d$ der Länge $k\ge1$ gilt
\[
 0\le V_b(d)<b^k.
\]
Für eine normalisierte Darstellung einer Zahl $n>0$ gilt genauer
\[
 b^{k-1}\le n<b^k.
\]
Diese beiden Ungleichungen charakterisieren ihre Stellenzahl $k$.
\end{theorem}

\section{Ziffern auswerten und zurückgewinnen}

\begin{theorem}[Ziffern durch Quotienten und Reste]
Ist $n=V_b(d)$ eine Darstellung der Länge $k\ge1$, so gilt
für $0\le i<k$
\[
 d_i=\bigl(n\mathbin{\mathrm{div}}b^i\bigr)\mathbin{\mathrm{mod}}b.
\]
Für jedes $j\in\mathbb N$ mit $j\le k$ gilt außerdem
\[
 n\mathbin{\mathrm{mod}}b^j=\sum_{i<j}d_i b^i,
 \qquad
 n\mathbin{\mathrm{div}}b^k=0.
\]
Bei $j=0$ sind beide Seiten der ersten Gleichung $0$.
\end{theorem}

\begin{theorem}[Horner-Auswertung]
Für Ziffern $d_0,\ldots,d_{k-1}$ mit $k\ge1$ seien
\[
 h_0=0,\qquad h_{j+1}=b h_j+d_{k-1-j}\quad(0\le j<k).
\]
Dann ist $h_k=V_b(d)$.
\end{theorem}

\begin{theorem}[Führende Nullen, Stellenverschiebung und Aneinanderfügen]
Eine führende Null ändert den Wert eines Zahlzeichens nicht.
Das Anhängen von $r$ Nullen rechts multipliziert seinen Wert mit $b^r$.
Für Ziffernblöcke $d$ der Länge $k\ge1$ und $e$ der Länge
$\ell\ge1$ gilt
\[
 V_b(d_{k-1}\cdots d_0e_{\ell-1}\cdots e_0)
 =b^\ell V_b(d)+V_b(e).
\]
Hier bezeichnet $V_b$ auf einem ausgeschriebenen Ziffernblock
dieselbe Stellenwertsumme wie in Definition~\ref{b10guide:def:digits}.
\end{theorem}

\begin{theorem}[Zahlenvergleich an der höchsten abweichenden Stelle]
Haben zwei verschiedene Ziffernbelegungen $d,e$ dieselbe Länge $k$,
und ist $j<k$ die höchste Stelle mit $d_j\ne e_j$, so gilt
\[
 V_b(d)<V_b(e)\quad\Longleftrightarrow\quad d_j<e_j.
\]
Bei normalisierten Darstellungen positiver Zahlen ist die Zahl
mit weniger Stellen stets die kleinere.
\end{theorem}

\section{Teilbarkeitsregeln und gebräuchliche Basen}

\begin{theorem}[Endziffern- und Quersummenregel]
Für $n=V_b(d)$ und $q\in\mathbb N_+$ gilt:
\[
\begin{aligned}
 q\mid b\quad&\Longrightarrow\quad
       \bigl(q\mid n\ \Longleftrightarrow\ q\mid d_0\bigr),\\
 q\mid(b-1)\quad&\Longrightarrow\quad
       \Bigl(q\mid n\ \Longleftrightarrow\ q\mid\sum_{i<k}d_i\Bigr).
\end{aligned}
\]
Insbesondere gilt $b\mid n\Longleftrightarrow d_0=0$.
Allgemeiner hängt für $1\le j\le k$ und $q\mid b^j$
die Teilbarkeit durch $q$ nur von den letzten $j$ Ziffern ab:
\[
 q\mid n\quad\Longleftrightarrow\quad q\mid\sum_{i<j}d_i b^i.
\]
\end{theorem}

\begin{definition}[Dualsystem und Dezimalsystem]
Das \emph{Dualsystem} ist das Stellenwertsystem zur Basis $2$ mit
Ziffern $0,1$; das \emph{Dezimalsystem} ist das Stellenwertsystem
zur Basis $10$ mit Ziffern $0,1,\ldots,9$.
Die Basisangabe kann im Dezimalsystem entfallen.
\end{definition}

\begin{theorem}[Teilbarkeitsregeln im Dezimalsystem]
Eine dezimal dargestellte natürliche Zahl ist genau dann
\begin{itemize}
 \item durch $2$ beziehungsweise $5$ teilbar, wenn ihre letzte
       Ziffer durch $2$ beziehungsweise $5$ teilbar ist;
 \item durch $4$ beziehungsweise $8$ teilbar, wenn der Wert ihrer
       letzten zwei beziehungsweise drei Ziffern entsprechend teilbar ist;
 \item durch $3$ beziehungsweise $9$ teilbar, wenn ihre Quersumme
       entsprechend teilbar ist.
\end{itemize}
Bei kürzeren Zahlzeichen werden für die Regel über zwei oder drei
Endziffern links Nullen ergänzt.
\end{theorem}

\begin{theorem}[Quersummen an geraden und ungeraden Stellen]
\label{b10guide:thm:alternating-digit-sums}
Sei $n=\sum_{i<k}d_i b^i$ eine Zifferndarstellung zur Basis $b\ge2$.
Die Summen der Ziffern an geraden beziehungsweise ungeraden Stellen seien
\[
 E:=\sum_{\substack{i<k\\i\text{ gerade}}}d_i,
 \qquad O:=\sum_{\substack{i<k\\i\text{ ungerade}}}d_i;
\]
eine leere Summe ist $0$. Für jedes $q\in\mathbb N_+$ mit $q\mid(b+1)$ gilt
\[
 q\mid n\quad\Longleftrightarrow\quad
 E\mathbin{\mathrm{mod}}q=O\mathbin{\mathrm{mod}}q.
\]
Insbesondere ist eine Dezimalzahl genau dann durch $11$ teilbar,
wenn ihre beiden Ziffernsummen denselben Rest bei Division durch $11$ haben.
Die Stellenzählung beginnt rechts mit $0$.
\end{theorem}

\section{Binomialkoeffizienten aus Ziffern}

\begin{theorem}[Satz von Lucas]
\label{b10guide:thm:lucas}
Seien $p$ eine Primzahl, $n,k\in\mathbb N$ und $\ell\ge1$ mit
$n<p^\ell$ und $k<p^\ell$. Schreibe, gegebenenfalls mit führenden Nullen,
\[
 n=\sum_{i<\ell}n_i p^i,\qquad
 k=\sum_{i<\ell}k_i p^i,\qquad n_i,k_i<p.
\]
Dann gilt
\[
 \binom{n}{k}\mathbin{\mathrm{mod}}p
 =\left(\prod_{i<\ell}\binom{n_i}{k_i}\right)
       \mathbin{\mathrm{mod}}p.
\]
Insbesondere charakterisieren die Ziffern die Nichtteilbarkeit:
\[
 p\nmid\binom{n}{k}\quad\Longleftrightarrow\quad
       \forall i<\ell\;(k_i\le n_i).
\]
Auch hier gilt $\binom{n}{k}=0$ für $k>n$; zusätzliche führende
Nullen ändern die Aussagen nicht \cite{natlucas}.
\end{theorem}

\Needspace{28\baselineskip}
\begin{theorem}[Kummers Übertragssatz und Quersummenformeln]
\label{b10guide:thm:kummer}
Seien $p$ eine Primzahl und $k,n\in\mathbb N$ mit $k\le n$.
Für $j<n$ setze
\[
 c_j=\begin{cases}
 1,&p^{j+1}\le
    (k\mathbin{\mathrm{mod}}p^{j+1})
    +((n-k)\mathbin{\mathrm{mod}}p^{j+1}),\\
 0,&\text{sonst}.
 \end{cases}
\]
Dann gilt
\[
 v_p\!\binom{n}{k}=\sum_{j<n}c_j.
\]
Dabei ist $c_j=1$ genau dann, wenn bei der Addition von $k$ und
$n-k$ zur Basis $p$ ein Übertrag aus den untersten $j+1$ Stellen
in die nächste Stelle erfolgt.

Für $t\in\mathbb N$ bezeichne $s_p(t)$ die Summe der Ziffern
seiner normalisierten Darstellung zur Basis $p$, also
$s_p(t)=\sum_{i<\ell}t_i$ für $t=\sum_{i<\ell}t_i p^i$;
insbesondere ist $s_p(0)=0$. Dann gelten die natürlichen Gleichungen
\[
 \begin{aligned}
 s_p(k)+s_p(n-k)
    &=s_p(n)+(p-1)v_p\!\binom{n}{k},\\
 n&=s_p(n)+(p-1)v_p(n!).
 \end{aligned}
\]
Die zweite Gleichung ist die Quersummenform der Legendre-Formel.
Für $n=0$ ist $k=0$, die Übertragssumme ist leer und die beiden
auftretenden Primexponenten sind $0$. Wegen $k\le n$ ist $\binom{n}{k}>0$;
ein Primexponent der Null wird nicht verwendet \cite{natlegendre}.
\end{theorem}
